%!TEX root = ../main.tex

\chapter{Methodology, Research Design, and Software Development}\label{ch:methodology}

  \section{Overall Research Design and Workflow}\label{sec:research-design}
  % Describe the overall research approach and data flow.
  % Explain the relationship between the chemistry and computational components.
  % Introduce the research datasets and evaluation strategy.

  \section{Environmental Samples and Chemistry Experimental Methods}\label{sec:environmental-samples}
  % Describe sample origins, locations, and sample types.
  % Document instrumentation and analytical conditions.
  % Describe LC-HRMS and MS/MS acquisition.
  % Explain positive- and negative-ion modes, where applicable.
  % Clarify which experimental work was performed by previous researchers.

  \section{Upstream Data Processing and Chemical Annotation}\label{sec:data-processing}
  % Document MS-DIAL processing and version.
  % Describe feature detection, alignment, and spectral exports, as applicable.
  % Explain data-quality checks and inclusion criteria.
  % Describe suspect screening, candidate annotations, and metadata.
  % Discuss identification confidence and limitations.

  \section{Software Requirements and System Architecture}\label{sec:software-architecture}
  % Define functional and nonfunctional requirements.
  % Document supported input formats and validation.
  % Describe system modules and data flow.
  % Explain data structures and metadata representation.
  % Describe reproducibility and provenance.

  \section{Spectral Comparison and Similarity Calculation}\label{sec:spectral-similarity}
  % Describe spectrum representation and preprocessing.
  % Explain peak matching and mass tolerances.
  % Document the cosine similarity implementation.
  % Define edge thresholds and computational considerations.
  % Explain handling of missing, incomplete, or unusable spectra.

  \section{Molecular-Network Construction}\label{sec:network-construction}
  % Describe node and edge representation.
  % Define similarity thresholds and graph filtering.
  % Explain connected components and isolated features.
  % Describe how network elements link to chemical metadata.

  \section{User Interface Design and Interaction}\label{sec:interface-design}
  % Describe network layout and visual encodings.
  % Explain node and edge selection.
  % Document feature metadata and similarity inspection.
  % Describe filtering, navigation, and network exploration.
  % Document saving, exporting, and persistence, where implemented.
  % Explain design decisions informed by chemist workflows.

  \section{Software Implementation and Reproducibility}\label{sec:implementation}
  % Document programming languages, libraries, and dependencies.
  % Describe modular design, version control, and documentation.
  % Explain error handling and logging.
  % Describe configuration, parameter recording, and reproducible execution.

  \section{Software Evaluation and Testing Methods}\label{sec:evaluation-methods}
  % Describe unit, integration, and end-to-end tests.
  % Explain reference calculations for similarity validation.
  % Document graph-construction correctness checks.
  % Describe repeated-run reproducibility tests.
  % Define runtime and scalability measurements.
  % Explain interface and workflow task evaluation.

  \section{Chemical Case-Study Design}\label{sec:case-study-design}
  % Explain the selection of representative environmental samples.
  % Define criteria for selecting samples, features, and network clusters.
  % Describe the use of existing annotations and spectral evidence.
  % Explain interpretation of known and unidentified features.
  % Define criteria for limiting chemical claims.
